% CPSY 1291 -- Recitation handout 7
% Compile:  latexmk -lualatex handout-07-midterm-review.tex
\documentclass[11pt]{article}
\usepackage{cpsy1291-handout}

\begin{document}

\handouthead{7}{Exam 1 review}
  {Lectures 1--8 and the four steps of backpropagation, in sixteen problems}
  {Week 7 \ $\cdot$\ Exam 1 is in class, Thu 10/22}

\section*{How to use this}
\addcontentsline{toc}{section}{How to use this}

This is a review, not new material. Everything here has appeared in a lecture,
an assignment, or an earlier recitation. Exam 1 covers Lectures 1--8 and the
four steps of backpropagation from Lecture 9; nothing else from Lecture 9 and
nothing after it.

Work in this order, and do not reverse it. \textbf{First} read Section~1, which
is the whole of that material compressed into the statements you should be able
to make without notes. \textbf{Then} do the problems in Section~2 with the
summary closed. Reading a summary produces the feeling of understanding;
answering a question produces the fact of it, and the gap between those two is
where exam marks are lost.

\begin{keybox}
\ask{What the exam can ask.} Every lecture in this course does three things:
introduces an \emph{object} (a distance, a matrix, a score, a rule), states
\emph{what it is for}, and states \emph{when it misleads you}. Exam questions
come from all three, and the third is the one students prepare least and are
asked about most.
\end{keybox}

\section{Lectures 1--8, compressed}

\subsection{L01 --- Representations, spaces and metrics}

A representation is a pattern that carries information about a stimulus.
Measured, it is a vector $\mathbf{x}\in\mathbb{R}^{D}$: 784 pixel values, 210
voxels, 482 IT firing rates, 4,096 network activations. $D$ is the number of
measurements, not the number of objects, and each image is one point in that
$D$-dimensional space. Stack $N$ images as rows and you have the
\textbf{response matrix} $\mathbf{X}\in\mathbb{R}^{N\times D}$, whose entry
$x_{ij}$ is feature $j$ of image $i$.

Single trials are noisy: two presentations of one image correlate at about 0.17
in monkey IT, and the median is 0.14 for fMRI. Averaging repeats keeps what
repeats, so every response vector in this course is a trial mean. The
\textbf{split-half correlation} --- the mean of half the repeats against the mean
of the other half (0.69 in the IT data) --- says how reliable that mean is.

Pixels, voxels, neurons and units cannot be compared entry by entry: the
vectors differ in length and in units. Their pairwise dissimilarities can. The
\textbf{RDM} $\mathbf{D}\in\mathbb{R}^{N\times N}$ holds
$d_{ik} = d(\mathbf{x}^{(i)},\mathbf{x}^{(k)})$, with a zero diagonal and
$d_{ik}=d_{ki}$. Choosing $d$ is a scientific decision.

\subsection{L02 --- Similarity, MDS and comparing RDMs}

\begin{center}
\begin{tabular}{@{}lll@{}}
\toprule
Measure & Computes & Ignores \\
\midrule
Euclidean & length of the difference vector & nothing: keeps gain and offset \\
cosine & $1-\cos\phi$ & gain (overall scale) \\
correlation & $1-r$: cosine on centered vectors & gain \emph{and} offset \\
\bottomrule
\end{tabular}
\end{center}

Choose by what should not count. With two components, $r=\pm1$ always, so
correlation needs more than two. A distance in the full sense (a metric)
satisfies nonnegativity, identity, symmetry and the triangle inequality.
Euclidean distance does; cosine and correlation fail identity (a brighter copy of
an image sits at cosine distance 0 from it) and the triangle inequality.

\textbf{MDS}: given $\mathbf{D}$, find coordinates $\mathbf{Y}\in\mathbb{R}^{N\times m}$
whose distances $\hat d_{ik}$ reproduce $d_{ik}$, by minimizing the stress
$\sum_{i<k}(\hat d_{ik}-d_{ik})^2$. Report \textbf{Stress-1}, which does not
depend on the units (Kruskal: 0.05 excellent, 0.10 fair, 0.20 poor), and compare
it with the stress of a shuffled matrix. Three warnings: the solution is not
unique (shift, rotate or reflect it and nothing changes, so the axes mean
nothing); $m$ is your choice (look for the elbow of stress against $m$); low
stress is not truth (enough dimensions fit noise). \textbf{Non-metric MDS} fits
only the rank order of the dissimilarities, which suits ratings: on the
Assignment~1 ratings Stress-1 falls from 0.31 to 0.23.

\textbf{Shepard's law}: in the psychological space recovered by non-metric MDS,
generalization falls exponentially with distance, $s = e^{-c\,d}$. The law is
not circular, because only the order was fitted. The fair rival is the Gaussian
$a\,e^{-bd^2}$, with the same number of parameters but flat at $d=0$, so the
test is at the origin. Against a network's distances the answer changes: that
is a different ruler.

\textbf{Comparing two RDMs} (RSA): flatten each upper triangle,
\texttt{np.triu\_indices(N, k=1)}, $N(N-1)/2$ numbers, and take Spearman's
$\rho$, which asks only whether the two systems order the pairs alike. $\rho$
alone means little: it needs a ceiling and a second model. Matching geometry is
evidence about a representation, not a demonstration of shared mechanism.

\subsection{L03 --- PCA and eigendecomposition}

Center each feature. The covariance matrix $\mathbf{C}$ has the variances on its
diagonal and summarizes the spread and orientation of the points. Project a
centered point onto a unit direction $\mathbf{w}$: $z=\mathbf{w}^\top\tilde{\mathbf{x}}$,
and across stimuli $\operatorname{Var}(z)=\mathbf{w}^\top\mathbf{C}\mathbf{w}$.
Projected variance and squared reconstruction error add to a fixed total, so
maximizing one minimizes the other.

Most directions are turned by $\mathbf{C}$. An \textbf{eigenvector} is not:
$\mathbf{C}\mathbf{v}_j=\lambda_j\mathbf{v}_j$, and its eigenvalue $\lambda_j$ is
the variance along it. PCA is the eigendecomposition
$\mathbf{C}=\mathbf{V}\boldsymbol{\Lambda}\mathbf{V}^\top$; PC1 is the eigenvector
with the largest eigenvalue. The \textbf{scores}
$z_j=\mathbf{v}_j^\top(\mathbf{x}-\boldsymbol{\mu})$ are the new coordinates, and
$\hat{\mathbf{x}}=\boldsymbol{\mu}+\sum_{j\le k}z_j\mathbf{v}_j$ is the
reconstruction. The scree plot and the cumulative curve give the
\textbf{embedding dimension} at a chosen threshold; a variance threshold is a
reconstruction budget, not a statement about any task. \texttt{PCA.fit\_transform}
centers each feature but does not z-score it.

Three cautions. \textbf{Scale}: with one feature in millimetres and one in
metres, PC1 is $1.00\,x_1+0.00\,x_2$; after z-scoring it is
$0.71\,x_1+0.71\,x_2$. Standardizing is a modelling decision, not a default.
\textbf{Most variance is not most useful}: the largest direction in the 120
photographs (29\%) is mostly brightness, and in the toy example PC2 carries 2.2\%
of the variance yet is the only direction that separates the groups.
\textbf{PCA is flat}: a curved path needs more linear coordinates than it has
degrees of freedom. In Bao et al., PCA of a network's responses gave an object
space whose directions are hypotheses, tested on held-out images.

\subsection{L04 --- Nonlinear dimensionality reduction}

A \textbf{manifold} is a curved surface that is flat up close. Three counts of
dimension answer three questions: \textbf{ambient} (coordinates per point),
\textbf{embedding} (the flat subspace PCA needs at a variance threshold) and
\textbf{intrinsic} (dimensions needed to locate a point on the manifold). For
the Swiss roll they are 3, 3 and 2. PCA overestimates the intrinsic dimension
whenever the manifold curves, and PCA and MDS both fold the roll, because points
on neighboring turns are close through the air.

\textbf{Isomap} measures along the sheet --- shortest paths through a
nearest-neighbor graph --- then runs MDS. One edge between turns, or a sparsely
sampled sheet, breaks it.

\textbf{$t$-SNE}: (1) in the data, a Gaussian around each point gives neighbor
probabilities $P$, its width set by the \textbf{perplexity}, roughly how many
neighbors each point listens to (typically 5--50, always below $N$); (2) in the
map, a heavy-tailed Student-$t$ gives $Q$, which leaves room for crowded
neighbors; (3) gradient descent on the KL divergence moves the map until $Q$
matches $P$. Close pairs dominate the mismatch, so $t$-SNE keeps neighbors, not
distances: cluster sizes and gaps are not measurements, low perplexity makes
clumps from noise, and another seed gives another map (PCA initialization is
more repeatable). \textbf{UMAP} builds a neighbor graph too
(\texttt{n\_neighbors}, \texttt{min\_dist}) and carries the same warnings.
Sweep the settings and the seed, keep what holds, and remember that stability is
not fidelity: a map can repeat itself and still disagree with the original
space. A map suggests a question; the response space answers it.

\subsection{L05 --- Geometry of neural representations}

\textbf{Rank} is the number of directions with nonzero variance. Noise makes
every real recording full rank (480 IT neurons: rank 480), so look at how the
variance spreads instead. Mouse V1 has no elbow: 90\% of the variance needs 744
components, and the spectrum is a power law with slope near 1. The
\textbf{participation ratio}
\[
\mathrm{PR}=\frac{\left(\sum_j\lambda_j\right)^2}{\sum_j\lambda_j^2},\qquad 1\le \mathrm{PR}\le D,
\]
needs no threshold: if $k$ directions share the variance equally, PR $=k$. It
counts how evenly the variance is spread, not how much information there is.
z-scoring each unit changes it (32 raw against 401 z-scored for one simulated
population), so report which you used.

A \textbf{linear readout} answers ``yes'' when $\mathbf{w}^\top\mathbf{x}+b>0$;
its boundary is a hyperplane, perpendicular to $\mathbf{w}$. In IT, animacy is
read out for 93\% of held-out views, against 50\% for guessing and 63\% for the
\textbf{majority-class baseline}. Two classes are \textbf{linearly separable}
when one hyperplane puts them on opposite sides. XOR is not (75\% at best); a
third unit $x_1x_2$ lifts the points so that one plane separates them.

\subsection{L06 --- Reading and comparing representations}

Two properties pull against each other. A representation is \textbf{flexible}
if almost any grouping of the stimuli can be read out --- measured by
\textbf{shattering} --- and \textbf{abstract} if a readout learned on some
stimuli still works on stimuli it never saw --- measured by \textbf{CCGP}.

$N$ stimuli in general position in $N-1$ dimensions can be split every way:
three points in two units give 8 of 8 splits, three on a line 6 of 8, four at
the corners of a square 14 of 16 (the XOR splits fail). Capacity is measured
with \textbf{random labels}: a linear readout memorizes about two per dimension,
and deep networks memorize ImageNet with shuffled labels (Zhang et al.\ 2017)
while staying at chance on new photographs.

A unit whose responses are a weighted sum of the others never raises the rank.
A nonlinear unit --- the product $x_1x_2$, or an AND --- does, and then all 16
splits are separable. That is why \textbf{nonlinear mixed selectivity}, common in
prefrontal cortex, makes a population high-dimensional: its shattering
dimensionality is close to the maximum of 24, against at most 8 for pure
selectivity, and it drops on error trials.

\textbf{CCGP}: fit a readout on some conditions (dogs, chairs) and test it,
weights fixed, on others (cats, tables). It is high when the new stimuli differ
from the old only along directions the readout ignores. Of Bernardi et al.'s
four toy geometries, points at random are flexible but not abstract, clusters
abstract but not flexible, a square abstract but not fully flexible, and a bent
square both. Hippocampus looks like the bent square: context CCGP 0.96 with high
shattering.

Comparing two systems: \textbf{RSA} builds an RDM per system and compares rank
orders (Spearman); \textbf{CKA} builds a Gram matrix (dot products of centered
units) and compares values, and is unchanged by uniform rescaling or rotation of
the units. CKA is more common for network layers, RSA for brains and judgments.

\subsection{L07 --- Geometry, baselines and tests}

On error trials, context CCGP drops but decoding of context (all conditions in
training) does not: the geometry breaks, the information does not.

A score means something only after three questions. \textbf{Is it above
chance?} Unrelated data correlate a little by luck. A \textbf{permutation test}
shuffles one system's image order (or a readout's labels) 1,000 times and asks
whether the score beats the 95th percentile of the shuffled scores (IT against
ResNet-18: 0.71 against 0.055). With more neurons than training images, any
labels fit (480 neurons, 102 images: 1.00 on training, chance on test), and with
few stimuli even unrelated systems reach a linear CKA near 1. \textbf{Would it
hold on new images?} Cross-validation tests every image once with a fit that
never saw it; \textbf{leakage} --- letting test images into any step, from the
z-scoring mean to the choice of units --- makes the score optimistic, so split
first and fit every step on the training images. A bootstrap over images says
whether a ranking of models would survive another sample of images.
\textbf{How high could it be?} Noise caps every score; the \textbf{noise
ceiling} comes from split-half reliability, and a score is read against it (0.3
is near the top if the ceiling is 0.35), alongside the shuffled floor and the
same area in another brain. Finally, change the measure and the winner can
change: RSA puts CLIP first, CKA ResNet-18.

\subsection{L08 --- Learning the weights}

The model neuron keeps what a point neuron does and drops the rest: inputs
$x_i$ (firing rates), weights $w_i$ (synaptic strengths, excitatory or
inhibitory), the pre-activation $z=\mathbf{w}^\top\mathbf{x}+b$, and the output
$a=g(z)$, most often a ReLU. Because
$\mathbf{w}^\top\mathbf{x}=\|\mathbf{w}\|\|\mathbf{x}\|\cos\phi$, the neuron
asks how much of $\mathbf{w}$ is in the input. Read as a \textbf{boundary},
$\mathbf{w}$ is perpendicular to the line $z=0$ and $b$ slides it: one neuron is
a linear classifier, the same readout as in L05--L07. Read as a
\textbf{template}, $\mathbf{w}$ drawn as an image looks like what the neuron
detects.

Learning needs a \textbf{loss} and \textbf{gradient descent},
$\mathbf{w}\leftarrow\mathbf{w}-\eta\nabla\ell$: the gradient points uphill, and
the learning rate $\eta$ sets the step.

\begin{center}
\begin{tabular}{@{}llll@{}}
\toprule
Rule & Output $\hat y$ & Loss & Update \\
\midrule
Perceptron (1958) & threshold & none: a heuristic & moves $\mathbf{w}$ only on mistakes \\
Delta rule (1960) & $\mathbf{w}^\top\mathbf{x}$ & $\tfrac12(y-\hat y)^2$ & $\Delta\mathbf{w}=\eta\,(y-\hat y)\,\mathbf{x}$ \\
Logistic regression & $\sigma(z)$ & cross-entropy & $\Delta\mathbf{w}=\eta\,(y-\hat p)\,\mathbf{x}$ \\
\bottomrule
\end{tabular}
\end{center}

The perceptron is the delta rule with a thresholded output, so its error is
zero unless it is wrong; if a line separates the classes, it stops. The
cross-entropy is 0 when the model is certain and right, 0.69 for a 50/50 guess,
and 2.30 when it gives the correct class 0.1. One recipe throughout:
\textbf{learning rate $\times$ error $\times$ input}.

\subsection{From L09 --- the four steps of backpropagation}

A hidden unit has no target, so its error has to come from the output's. With
a sigmoid output and the cross-entropy:

\begin{enumerate}
  \item \textbf{Forward}: run the input through and keep every weighted sum $z$
        and every output $a$.
  \item \textbf{Output error}: $\delta_{\text{out}}=\partial\ell/\partial z_{\text{out}}=\hat y-y$.
  \item \textbf{Backward}: each hidden unit's error is the errors of the units it
        feeds, each times the weight to them, all times its own slope,
        $\delta_j=g'(z_j)\sum_k w_{kj}\,\delta_k$. A ReLU's slope is 1 if it was
        on and 0 if it was off, so a unit that was off gets no error.
  \item \textbf{Gradients}: each weight's gradient is the error of the unit it
        enters times its input, $\partial\ell/\partial w=\delta\times\text{input}$;
        each bias gets its unit's error.
\end{enumerate}

The lecture's 2--2--1 XOR network, input $(1,0)$, target 1: $z^{[1]}=(0.50,0.30)$,
$z^{[2]}=0.17$, $\hat y=0.54$, loss 0.61; $\delta^{[2]}=-0.46$;
$\delta^{[1]}=(-0.32,\,0.27)$; $\partial\ell/\partial w^{[2]}_1=-0.46\times0.50=-0.23$.
One step with $\eta=1$ raises $\hat y$ to 0.82.

\section{Problems}

Do these closed-book. Answers follow.

\ask{1.} A unit's activity is $(2, 4, 6)$ across three stimuli; another's is
$(1, 2, 3)$. Give the cosine distance and the correlation distance between them,
and explain the pair of answers in one sentence.

\ask{2.} Why is correlation distance undefined for a unit that produces the same
value for every stimulus?

\ask{3.} You are told $d(A,B) = 0.2$, $d(B,C) = 0.2$, $d(A,C) = 0.9$. Which
requirement of a distance fails, and name a measure from this course that
permits it.

\ask{4.} Someone reports that an MDS map's horizontal axis is ``animacy''. What
is the methodological objection, and what evidence would answer it?

\ask{5.} Two features are recorded, one in millimetres and one in metres. What
is PC1 likely to be in raw units, what changes after z-scoring, and when should
you not z-score?

\ask{6.} PC1 keeps 97.8\% of the variance and PC2 2.2\%. Should you drop PC2
before training a readout to separate two groups?

\ask{7.} Give the ambient, embedding and intrinsic dimension of a Swiss roll
(a 2-D sheet rolled up in 3-D), and say which of them PCA measures.

\ask{8.} Two $t$-SNE runs on identical data, differing only in
\texttt{random\_state}, give visibly different maps. Is either wrong? What
should be reported, and how would you check a grouping?

\ask{9.} Four units have covariance eigenvalues $(4, 4, 0, 0)$; another four have
$(7, 1, 0, 0)$. Give the rank and the participation ratio of each, and say what
the difference means.

\ask{10.} Four stimuli sit at the corners of a square in two units. How many of
the 16 splits can a linear readout separate, which fail, and what kind of third
unit fixes them?

\ask{11.} On error trials, context CCGP drops but decoding of context does not.
What has changed in the representation, and what has not?

\ask{12.} A readout on 480 IT neurons, trained on 102 images with shuffled labels,
scores 1.00 on its training images. Explain, and say what you report instead.

\ask{13.} You z-score your responses using the mean and standard deviation of all
images, then run cross-validation. What is wrong, and how large is the problem
compared with choosing units by their test score?

\ask{14.} Delta rule: $\mathbf{w}=(0.5,-1)$, $b=0$, $\mathbf{x}=(2,1)$, $y=1$,
$\eta=0.1$. Compute $\hat y$, $\Delta\mathbf{w}$, and $\hat y$ after the update.

\ask{15.} A logistic-regression neuron gives $\hat p=0.5$ for an example of class
1, and another gives $\hat p=0.1$. Give each cross-entropy, and the update
$\Delta\mathbf{w}$ in terms of $\mathbf{x}$.

\ask{16.} The lecture's XOR network with its starting weights (hidden
$\mathbf{w}_1=(0.50,-0.40)$, $\mathbf{w}_2=(0.30,0.80)$, output weights
$(0.70,-0.60)$, biases 0, ReLU hidden units, sigmoid output, cross-entropy).
Input $(0,1)$, target 1. Run the four steps of backpropagation and give every
gradient.

\subsection*{Answers}

\ask{1.} Cosine distance $0$: the two vectors point in exactly the same
direction, differing only by a factor of two. Correlation distance $0$ as well:
after centering, both become proportional to $(-1, 0, 1)$. The units differ only
in gain, which cosine discards and correlation discards along with offset, so
neither can tell them apart; Euclidean distance keeps gain and can.

\ask{2.} Centering a constant vector gives the zero vector, whose length is zero,
and the correlation divides by that length. In code it surfaces as \texttt{nan},
which then propagates silently through the whole matrix.

\ask{3.} The triangle inequality: $0.9 > 0.2 + 0.2$. Cosine distance and
correlation distance both permit it; Euclidean distance does not.

\ask{4.} The MDS solution can be shifted, rotated or reflected without changing
any distance or the stress, so ``the horizontal axis'' is not a property of the
data; only relative positions are. Answering it requires independent evidence:
relate a direction in the map to an external animacy rating, or show that a
readout for animacy works on held-out stimuli.

\ask{5.} PCA maximizes variance, so the feature in millimetres dominates: PC1 is
close to $1.00\,x_1+0.00\,x_2$ and the second feature is invisible. After
z-scoring, PC1 weights both (0.71 each in the lecture's example) and recovers
the structure they share. Do not z-score by default when the units are
comparable and the variances differ for real reasons: standardizing changes the
relative weighting of the features.

\ask{6.} No. PCA ranks directions by variance; the task decides what is useful.
In the lecture's example both groups have the same PC1 coordinates, and a
threshold at zero on PC2 separates them.

\ask{7.} Ambient 3, embedding 3, intrinsic 2. PCA measures the embedding
dimension, a linear count that needs a variance threshold, and overestimates the
intrinsic dimension whenever the manifold curves.

\ask{8.} Neither is wrong: the random start is set by the seed. Report what is
stable across seeds and across a perplexity sweep, and say how many runs you
made. Sizes and gaps in the map are not distances in the data; to check a
grouping, compare each image's nearest neighbors in the original response space,
or test a readout on held-out images.

\ask{9.} Both have rank 2. PR $=8^2/(16+16)=2$ for the first and
$8^2/(49+1)=1.28$ for the second. Both use two directions, but in the second one
direction carries almost all the variance; PR counts how evenly the variance is
spread, not how much information the population carries.

\ask{10.} 14 of 16. The two XOR splits, which put opposite corners together,
fail. A third unit whose responses are not a weighted sum of the first two
--- $x_1x_2$, or an AND unit --- raises the rank to 3, and all 16 become
separable. Adding $x_1+x_2$ leaves the rank at 2.

\ask{11.} The abstract format of context has broken: a readout trained on some
conditions no longer transfers to the others. The information has not gone:
a readout trained on all conditions still decodes context.

\ask{12.} With more neurons than training images, the points are nearly
perpendicular and any labelling can be fit; scoring 1.00 means memorizing.
Report accuracy on held-out images, against the shuffled-label score on test
images (0.56 in the lecture, chance) and the majority-class baseline.

\ask{13.} Leakage: the test images influenced the z-scoring, so the test score
is somewhat optimistic. It is a small leak. Choosing units or settings by their
test score is a large one --- even pure noise can then look predictive. Split
first, then fit every step on the training images.

\ask{14.} $\hat y=0.5\cdot2-1\cdot1=0$. Error $y-\hat y=1$, so
$\Delta\mathbf{w}=0.1\times1\times(2,1)=(0.2,0.1)$ and $\mathbf{w}$ becomes
$(0.7,-0.9)$. After the update $\hat y=1.4-0.9=0.5$, and the squared-error loss
falls from 0.5 to 0.125.

\ask{15.} $-\log0.5=0.69$ and $-\log0.1=2.30$. The update is
$\Delta\mathbf{w}=\eta\,(y-\hat p)\,\mathbf{x}$: $0.5\,\eta\,\mathbf{x}$ for the
first and $0.9\,\eta\,\mathbf{x}$ for the second, so the more confidently wrong
neuron moves more.

\ask{16.} \emph{Step 1.} $z^{[1]}_1=-0.40$, so $a^{[1]}_1=0$ (off);
$z^{[1]}_2=0.80$, so $a^{[1]}_2=0.80$. $z^{[2]}=0.70\cdot0-0.60\cdot0.80=-0.48$,
$\hat y=\sigma(-0.48)=0.38$, loss $-\log0.38=0.96$. \emph{Step 2.}
$\delta^{[2]}=\hat y-y=-0.62$. \emph{Step 3.}
$\delta^{[1]}_1=-0.62\times0.70\times0=0$: the unit was off and gets no error.
$\delta^{[1]}_2=-0.62\times(-0.60)\times1=0.37$. \emph{Step 4.} Output weights:
$-0.62\times0=0$ and $-0.62\times0.80=-0.49$. Unit~1's weights: $(0,0)$.
Unit~2's weights: $0.37\times(0,1)=(0,\,0.37)$. Biases: $-0.62$ (output), $0$
and $0.37$ (hidden). The ReLU gates the error: nothing into unit~1 moves on this
example. One step with $\eta=1$ raises $\hat y$ to 0.65.

\section*{The night before}
\addcontentsline{toc}{section}{The night before}

Six things worth being able to state cold:

\begin{enumerate}
  \item the three dissimilarity measures, what each ignores, and which are
        metrics;
  \item what an eigenvector of the covariance is, in a sentence with no algebra;
  \item ambient, embedding and intrinsic dimension, and what a $t$-SNE map cannot
        show;
  \item what shattering measures and what CCGP measures, and why they pull
        against each other;
  \item the three questions before trusting a score: above chance, held out,
        against its ceiling;
  \item one update for three rules (learning rate $\times$ error $\times$ input),
        and the four steps of backpropagation.
\end{enumerate}

\end{document}
